% Options for packages loaded elsewhere
% Options for packages loaded elsewhere
\PassOptionsToPackage{unicode}{hyperref}
\PassOptionsToPackage{hyphens}{url}
\PassOptionsToPackage{dvipsnames,svgnames,x11names}{xcolor}
%
\documentclass[
  letterpaper,
  DIV=11,
  numbers=noendperiod]{scrartcl}
\usepackage{xcolor}
\usepackage{amsmath,amssymb}
\setcounter{secnumdepth}{-\maxdimen} % remove section numbering
\usepackage{iftex}
\ifPDFTeX
  \usepackage[T1]{fontenc}
  \usepackage[utf8]{inputenc}
  \usepackage{textcomp} % provide euro and other symbols
\else % if luatex or xetex
  \usepackage{unicode-math} % this also loads fontspec
  \defaultfontfeatures{Scale=MatchLowercase}
  \defaultfontfeatures[\rmfamily]{Ligatures=TeX,Scale=1}
\fi
\usepackage{lmodern}
\ifPDFTeX\else
  % xetex/luatex font selection
\fi
% Use upquote if available, for straight quotes in verbatim environments
\IfFileExists{upquote.sty}{\usepackage{upquote}}{}
\IfFileExists{microtype.sty}{% use microtype if available
  \usepackage[]{microtype}
  \UseMicrotypeSet[protrusion]{basicmath} % disable protrusion for tt fonts
}{}
\makeatletter
\@ifundefined{KOMAClassName}{% if non-KOMA class
  \IfFileExists{parskip.sty}{%
    \usepackage{parskip}
  }{% else
    \setlength{\parindent}{0pt}
    \setlength{\parskip}{6pt plus 2pt minus 1pt}}
}{% if KOMA class
  \KOMAoptions{parskip=half}}
\makeatother
% Make \paragraph and \subparagraph free-standing
\makeatletter
\ifx\paragraph\undefined\else
  \let\oldparagraph\paragraph
  \renewcommand{\paragraph}{
    \@ifstar
      \xxxParagraphStar
      \xxxParagraphNoStar
  }
  \newcommand{\xxxParagraphStar}[1]{\oldparagraph*{#1}\mbox{}}
  \newcommand{\xxxParagraphNoStar}[1]{\oldparagraph{#1}\mbox{}}
\fi
\ifx\subparagraph\undefined\else
  \let\oldsubparagraph\subparagraph
  \renewcommand{\subparagraph}{
    \@ifstar
      \xxxSubParagraphStar
      \xxxSubParagraphNoStar
  }
  \newcommand{\xxxSubParagraphStar}[1]{\oldsubparagraph*{#1}\mbox{}}
  \newcommand{\xxxSubParagraphNoStar}[1]{\oldsubparagraph{#1}\mbox{}}
\fi
\makeatother


\usepackage{longtable,booktabs,array}
\usepackage{calc} % for calculating minipage widths
% Correct order of tables after \paragraph or \subparagraph
\usepackage{etoolbox}
\makeatletter
\patchcmd\longtable{\par}{\if@noskipsec\mbox{}\fi\par}{}{}
\makeatother
% Allow footnotes in longtable head/foot
\IfFileExists{footnotehyper.sty}{\usepackage{footnotehyper}}{\usepackage{footnote}}
\makesavenoteenv{longtable}
\usepackage{graphicx}
\makeatletter
\newsavebox\pandoc@box
\newcommand*\pandocbounded[1]{% scales image to fit in text height/width
  \sbox\pandoc@box{#1}%
  \Gscale@div\@tempa{\textheight}{\dimexpr\ht\pandoc@box+\dp\pandoc@box\relax}%
  \Gscale@div\@tempb{\linewidth}{\wd\pandoc@box}%
  \ifdim\@tempb\p@<\@tempa\p@\let\@tempa\@tempb\fi% select the smaller of both
  \ifdim\@tempa\p@<\p@\scalebox{\@tempa}{\usebox\pandoc@box}%
  \else\usebox{\pandoc@box}%
  \fi%
}
% Set default figure placement to htbp
\def\fps@figure{htbp}
\makeatother





\setlength{\emergencystretch}{3em} % prevent overfull lines

\providecommand{\tightlist}{%
  \setlength{\itemsep}{0pt}\setlength{\parskip}{0pt}}



 


\usepackage{graphicx}
\usepackage{grffile}
\KOMAoption{captions}{tableheading}
\makeatletter
\@ifpackageloaded{caption}{}{\usepackage{caption}}
\AtBeginDocument{%
\ifdefined\contentsname
  \renewcommand*\contentsname{Table of contents}
\else
  \newcommand\contentsname{Table of contents}
\fi
\ifdefined\listfigurename
  \renewcommand*\listfigurename{List of Figures}
\else
  \newcommand\listfigurename{List of Figures}
\fi
\ifdefined\listtablename
  \renewcommand*\listtablename{List of Tables}
\else
  \newcommand\listtablename{List of Tables}
\fi
\ifdefined\figurename
  \renewcommand*\figurename{Figure}
\else
  \newcommand\figurename{Figure}
\fi
\ifdefined\tablename
  \renewcommand*\tablename{Table}
\else
  \newcommand\tablename{Table}
\fi
}
\@ifpackageloaded{float}{}{\usepackage{float}}
\floatstyle{ruled}
\@ifundefined{c@chapter}{\newfloat{codelisting}{h}{lop}}{\newfloat{codelisting}{h}{lop}[chapter]}
\floatname{codelisting}{Listing}
\newcommand*\listoflistings{\listof{codelisting}{List of Listings}}
\makeatother
\makeatletter
\makeatother
\makeatletter
\@ifpackageloaded{caption}{}{\usepackage{caption}}
\@ifpackageloaded{subcaption}{}{\usepackage{subcaption}}
\makeatother
\usepackage{bookmark}
\IfFileExists{xurl.sty}{\usepackage{xurl}}{} % add URL line breaks if available
\urlstyle{same}
\hypersetup{
  colorlinks=true,
  linkcolor={blue},
  filecolor={Maroon},
  citecolor={Blue},
  urlcolor={Blue},
  pdfcreator={LaTeX via pandoc}}


\author{}
\date{}
\begin{document}


\section{Smale 的工作：h-Cobordism
定理与广义庞加莱猜想}\label{smale-ux7684ux5de5ux4f5ch-cobordism-ux5b9aux7406ux4e0eux5e7fux4e49ux5e9eux52a0ux83b1ux731cux60f3}

\begin{quote}
📚 \textbf{参考书目}：\emph{Introductory Lectures on Manifold Topology
Signposts} (Thomas Farrell \& Yang Su)\\
🏷️ \textbf{标签}：\texttt{\#smale} \texttt{\#h-cobordism}
\texttt{\#manifold}
\end{quote}

\begin{center}\rule{0.5\linewidth}{0.5pt}\end{center}

\section{1
拓扑广义庞加莱猜想}\label{ux62d3ux6251ux5e7fux4e49ux5e9eux52a0ux83b1ux731cux60f3}

\subsection{1.1 主要结果}\label{ux4e3bux8981ux7ed3ux679c}

\begin{quote}
{[}!important{]} 🎯 Smale 定理 (GPC: 拓扑广义庞加莱猜想)

设 \(M^m\) 是闭光滑流形，维数 \(m \geq 5\)。若 \(M\) 是同伦球面，则
\(M\) \textbf{同胚于球面}。

\textbf{主要工具} (\(m \geq 5\))：\textbf{h-Cobordism 定理}

\begin{quote}
设 \(W^m\) 是光滑的 h-cobordism，若 \(W\) 单连通且维数
\(m \geq 6\)，则配边 \(W\) 微分同胚于
\(\partial^- W \times [0,1]\)（相对于边界 \(\partial^- W\)）。
\end{quote}
\end{quote}

更进一步，利用 h-cobordism
定理，对于单连通光滑流形（特别是光滑同伦球面），当 \(n \geq 5\)
时有\textbf{双射}：

\[\boxed{\frac{\{S^n \text{ 上的光滑结构}\}}{\{微分同胚\}} \longleftrightarrow \frac{\{n \text{ 维同伦球面}\}}{\{\text{h-cobordism}\}}}\]

\begin{longtable}[]{@{}cl@{}}
\toprule\noalign{}
& 说明 \\
\midrule\noalign{}
\endhead
\bottomrule\noalign{}
\endlastfoot
\textbf{左边} & 球面上所有光滑结构，即所有怪球构成的集合 \\
\textbf{右边} & 所有同伦球面的配边类，记为 \(\Theta^n\) \\
\end{longtable}

\textbf{\(\Theta^n\) 的群结构}：

\begin{longtable}[]{@{}
  >{\centering\arraybackslash}p{(\linewidth - 2\tabcolsep) * \real{0.5556}}
  >{\raggedright\arraybackslash}p{(\linewidth - 2\tabcolsep) * \real{0.4444}}@{}}
\toprule\noalign{}
\begin{minipage}[b]{\linewidth}\centering
运算
\end{minipage} & \begin{minipage}[b]{\linewidth}\raggedright
定义
\end{minipage} \\
\midrule\noalign{}
\endhead
\bottomrule\noalign{}
\endlastfoot
乘法 & 连通和 \\
单位元 & \(S^n\) 上的标准光滑结构 \\
逆元 & 反转定向，记为 \(-\Sigma\)（在柱 \(\Sigma \times I\) 内挖去
\(D^{n+1}\)） \\
\end{longtable}

\textbf{\(\Theta^n\) 的阶数}（与伯努利数密切相关）：

\begin{figure}[H]

{\centering \pandocbounded{\includegraphics[keepaspectratio]{2-Smale_s_work.png}}

}

\caption{怪异球面个数表}

\end{figure}%

\begin{quote}
{[}!warning{]} ⚠️ 注意 由于 h-cobordism 定理在 4 维失效（h-cobordism
类中的元素未必微分同胚），\(\Theta^4 = 1\) 不能推出 4 维光滑庞加莱猜想。
\end{quote}

\begin{center}\rule{0.5\linewidth}{0.5pt}\end{center}

\subsection{1.2 GPC 的证明}\label{gpc-ux7684ux8bc1ux660e}

\emph{本节假设 h-cobordism 定理成立。}

\subsubsection{1.2.1 h-Cobordism
的定义}\label{h-cobordism-ux7684ux5b9aux4e49}

\begin{quote}
{[}!info{]} 📐 定义 (Cobordism)

\((m+1)\) 维 cobordism 是流形 \(W^{m+1}\)，其边界可写成不交并：
\[\partial W = \partial^- W \sqcup \partial^+ W\] 称 \(W\) 是从
\(\partial^- W\) 到 \(\partial^+ W\) 的\textbf{配边}。

\emph{例：配边但不是 h-配边} 👇

\begin{figure}[H]

{\centering \pandocbounded{\includegraphics[keepaspectratio]{3-Smale_s_work.png}}

}

\caption{配边示例}

\end{figure}%
\end{quote}

\begin{quote}
{[}!info{]} 📐 定义 (h-Cobordism)

\((m+1)\) 维 \textbf{h-cobordism} 是配边 \(W^{m+1}\)，满足含入映射
\[\partial^- W \hookrightarrow W, \quad \partial^+ W \hookrightarrow W\]
都是\textbf{同伦等价}。

\emph{例：平凡配边} 👇

\begin{figure}[H]

{\centering \pandocbounded{\includegraphics[keepaspectratio]{4-Smale_s_work.png}}

}

\caption{平凡配边}

\end{figure}%
\end{quote}

\subsubsection{1.2.2 h-Cobordism 定理}\label{h-cobordism-ux5b9aux7406}

\begin{quote}
{[}!important{]} 🎯 h-Cobordism 定理

设 \(W^{m+1}\) 是光滑的 h-cobordism。若 \(W\) \textbf{单连通}且
\(m \geq 5\)，则
\[W \cong_{\text{diff}} \partial^- W \times [0,1] \quad (\text{rel } \partial^- W)\]
\end{quote}

\subsubsection{\texorpdfstring{1.2.3 GPC 的证明
(\(n \geq 6\))}{1.2.3 GPC 的证明 (n \textbackslash geq 6)}}\label{gpc-ux7684ux8bc1ux660e-n-geq-6}

\begin{quote}
{[}!tip{]} ✏️ 证明

设 \(M^m\) 是光滑流形，维数 \(\geq 6\) 且是同伦球面。

\begin{center}\rule{0.5\linewidth}{0.5pt}\end{center}

\textbf{Step 1️⃣} 在 \(M^m\) 上挖去两个隔开的 disk，得到配边 \(W^m\)。

\begin{figure}[H]

{\centering \pandocbounded{\includegraphics[keepaspectratio]{5-Smale_s_work.png}}

}

\caption{挖去disk}

\end{figure}%

\begin{center}\rule{0.5\linewidth}{0.5pt}\end{center}

\textbf{Step 2️⃣} 由 \textbf{Lemma 2.6}，\(W\) 单连通且是光滑
h-cobordism。

\begin{center}\rule{0.5\linewidth}{0.5pt}\end{center}

\textbf{Step 3️⃣} 应用 h-cobordism 定理：
\[W \cong_{\text{diff}} S^{m-1} \times [0,1] \quad (\text{rel } \partial^- W = S^{m-1})\]

\begin{center}\rule{0.5\linewidth}{0.5pt}\end{center}

\textbf{Step 4️⃣} 重新粘回挖去的 disk： - 底部 disk：通过 identity
映射粘回 → 得到 \(\mathbb{D}^m\) - 顶部
disk：任意粘接映射（影响光滑结构，不影响拓扑）→ 仍同胚于 \(S^m\)

\begin{figure}[H]

{\centering \pandocbounded{\includegraphics[keepaspectratio]{1-Smale_s_work.png}}

}

\caption{同胚构造}

\end{figure}%

\begin{center}\rule{0.5\linewidth}{0.5pt}\end{center}

\[\boxed{M \approx S^m} \quad \blacksquare\]
\end{quote}

\begin{center}\rule{0.5\linewidth}{0.5pt}\end{center}

\begin{quote}
{[}!note{]} 📜 Lemma 2.6 同伦球面上挖去两个 disk 得到 h-cobordism
(\(n \geq 6\))

\textbf{设置}：从同伦球面 \(M\) 挖去 \(D^+, D^-\)，得配边 \(W\)。

\begin{figure}[H]

{\centering \pandocbounded{\includegraphics[keepaspectratio]{7-Smale_s_work.png}}

}

\caption{设置示意图}

\end{figure}%

\begin{center}\rule{0.5\linewidth}{0.5pt}\end{center}

\textbf{证明}：

\textbf{(I) \(W\) 单连通} --- 从「\(M\) 单连通」推出「\(W\) 单连通」

\begin{longtable}[]{@{}lcl@{}}
\toprule\noalign{}
Triple & 定理 & 结论 \\
\midrule\noalign{}
\endhead
\bottomrule\noalign{}
\endlastfoot
\((M, W \cup D^+, D^-)\) & Van Kampen & \(W \cup D^+\) 单连通 \\
\((W \cup D^+, W, D^+)\) & Van Kampen & \(W\) 单连通 ✓ \\
\end{longtable}

\begin{center}\rule{0.5\linewidth}{0.5pt}\end{center}

\textbf{(II) \(\partial^+ W \hookrightarrow W\) 是同伦等价}

\begin{enumerate}
\def\labelenumi{\arabic{enumi}.}
\tightlist
\item
  对 \((W \cup D^+, W, D^+)\) 用 \textbf{Mayer-Vietoris} → 同调群同构
\item
  用 \textbf{Homology Whitehead 定理}：单连通 + 同调等价 ⟹ 同伦等价
\end{enumerate}

\begin{figure}[H]

{\centering \pandocbounded{\includegraphics[keepaspectratio]{9-Smale_s_work.png}}

}

\caption{Whitehead定理}

\end{figure}%

\begin{center}\rule{0.5\linewidth}{0.5pt}\end{center}

\textbf{(III)} \(\partial^- W \hookrightarrow W\) 类似。\(\blacksquare\)
\end{quote}

\begin{center}\rule{0.5\linewidth}{0.5pt}\end{center}

\begin{quote}
{[}!abstract{]} 📎 附：Whitehead 定理

\begin{longtable}[]{@{}
  >{\raggedright\arraybackslash}p{(\linewidth - 2\tabcolsep) * \real{0.5000}}
  >{\raggedright\arraybackslash}p{(\linewidth - 2\tabcolsep) * \real{0.5000}}@{}}
\toprule\noalign{}
\begin{minipage}[b]{\linewidth}\raggedright
概念
\end{minipage} & \begin{minipage}[b]{\linewidth}\raggedright
定义
\end{minipage} \\
\midrule\noalign{}
\endhead
\bottomrule\noalign{}
\endlastfoot
\textbf{弱同伦等价} & 存在 \(f: X \to Y\) 使
\(f_*: \pi_i(X) \xrightarrow{\cong} \pi_i(Y)\) 对 \(\forall i\) \\
\textbf{Whitehead 定理} & CW 复形间：弱同伦等价 ⟹ 同伦等价 \\
\textbf{Homology Whitehead} & 单连通空间：同调等价 ⟹ 同伦等价 \\
\end{longtable}

💡 弱同伦等价要求存在映射 \(f\)
实现同构，而非仅同伦群同构。存在同伦群同构但不同伦等价的反例。
\end{quote}

\begin{center}\rule{0.5\linewidth}{0.5pt}\end{center}

\section{2 h-Cobordism 定理}\label{h-cobordism-ux5b9aux7406-1}

\begin{quote}
{[}!info{]} 📐 Isotopy（同痕）

\begin{figure}[H]

{\centering \pandocbounded{\includegraphics[keepaspectratio]{10-Smale_s_work.png}}

}

\caption{同痕定义}

\end{figure}%

我们讨论的同痕总是\textbf{光滑同痕}。

⚠️ 连续同痕不能区分纽结！

\begin{figure}[H]

{\centering \pandocbounded{\includegraphics[keepaspectratio]{11-Smale_s_work.png}}

}

\caption{纽结示例}

\end{figure}%
\end{quote}

\begin{center}\rule{0.5\linewidth}{0.5pt}\end{center}

\subsection{2.1 Handle 分解}\label{handle-ux5206ux89e3}

\emph{（粘接映射）}

\begin{center}\rule{0.5\linewidth}{0.5pt}\end{center}

\subsection{2.2 Handle 运算法则}\label{handle-ux8fd0ux7b97ux6cd5ux5219}

\subsubsection{2.2.1 同痕不变性}\label{ux540cux75d5ux4e0dux53d8ux6027}

\begin{quote}
{[}!note{]} 📜 Prop：同痕扩张定理

\begin{figure}[H]

{\centering \pandocbounded{\includegraphics[keepaspectratio]{12-Smale_s_work.png}}

}

\caption{同痕扩张定理}

\end{figure}%
\end{quote}

\begin{quote}
{[}!note{]} 📜 Lemma 1：同痕引理

将 handle \(h\) 粘到流形 \(M\) 上，若使用\textbf{同痕}的粘接映射
\(f, g: \partial^- h \to \partial M\)，则
\[\boxed{M +_f h \cong M +_g h}\]

\textbf{证明}：

\begin{figure}[H]

{\centering \pandocbounded{\includegraphics[keepaspectratio]{13-Smale_s_work.png}}

}

\caption{同痕引理证明}

\end{figure}%
\end{quote}

\begin{center}\rule{0.5\linewidth}{0.5pt}\end{center}

\subsubsection{2.2.2
粘接顺序的交换性}\label{ux7c98ux63a5ux987aux5e8fux7684ux4ea4ux6362ux6027}

\begin{quote}
{[}!info{]} 📐 定义：Attaching sphere 与 Complementary sphere

Handle 写成乘积 \(D^i \times D^{n-i}\)：

\begin{longtable}[]{@{}
  >{\centering\arraybackslash}p{(\linewidth - 6\tabcolsep) * \real{0.2500}}
  >{\centering\arraybackslash}p{(\linewidth - 6\tabcolsep) * \real{0.2500}}
  >{\centering\arraybackslash}p{(\linewidth - 6\tabcolsep) * \real{0.2500}}
  >{\centering\arraybackslash}p{(\linewidth - 6\tabcolsep) * \real{0.2500}}@{}}
\toprule\noalign{}
\begin{minipage}[b]{\linewidth}\centering
分量
\end{minipage} & \begin{minipage}[b]{\linewidth}\centering
名称
\end{minipage} & \begin{minipage}[b]{\linewidth}\centering
Disk
\end{minipage} & \begin{minipage}[b]{\linewidth}\centering
Sphere
\end{minipage} \\
\midrule\noalign{}
\endhead
\bottomrule\noalign{}
\endlastfoot
第一（粘接用） & Attaching part & \(D^i \times \{0\}\) & \(\mathscr{S}\)
(attaching sphere) \\
第二（补维数） & Complementary part & \(\{0\} \times D^{n-i}\) &
\(\mathcal{S}\) (complementary sphere) \\
\end{longtable}

\begin{figure}[H]

{\centering \pandocbounded{\includegraphics[keepaspectratio]{14-Smale_s_work.png}}

}

\caption{Attaching和Complementary}

\end{figure}%
\end{quote}

\begin{quote}
{[}!important{]} 🎯 Thom 横截性定理

对三元组 \((W, M, N)\)（\(M, N\) 是 \(W\) 的子流形）：

\begin{enumerate}
\def\labelenumi{\arabic{enumi}.}
\tightlist
\item
  可做微小扰动（同痕），使 \(M, N\) \textbf{横截相交}
\item
  横截相交时：\(\dim M + \dim N = \dim W + \dim(M \cap N)\)
\end{enumerate}

\begin{center}\rule{0.5\linewidth}{0.5pt}\end{center}

\textbf{🔍 例}：\(\dim W = 3\)，\(\dim M = \dim N = 1\)

\[\dim(M \cap N) = 1 + 1 - 3 = -1 \implies M \cap N = \varnothing\]

💡 \textbf{启示}：通过控制维数可保证子流形不相交！
\end{quote}

\begin{quote}
{[}!note{]} 📜 Lemma 2：顺序引理

考虑 \((M + h_1) + h_2\)。若 \(h_2\) 的粘接球面 \(\mathscr{S}\) 与
\(h_1\) 的补球面 \(\mathcal{S}\) \textbf{不相交}，则可交换顺序：
\[\boxed{(M + h_1) + h_2 \cong (M + h_2) + h_1}\] 记作
\(M + (h_1 \sqcup h_2)\)。

\textbf{证明}：不相交时，可在 \(\partial(M + h_1)\) 上挖去
\(\mathcal{S}\)，将 \(\partial^- h_2\) 收缩到 \(\partial M\)。

\begin{figure}[H]

{\centering \pandocbounded{\includegraphics[keepaspectratio]{17-Smale_s_work.png}}

}

\caption{顺序引理证明}

\end{figure}%
\end{quote}

\begin{quote}
{[}!note{]} 📜 Coro

若 \(\mathrm{index}(h_2) \leq \mathrm{index}(h_1)\)，则
\[M + h_1 + h_2 = M + h_2 + h_1\]

💡 即：总可从 \textbf{index 更小}的 handle 开始粘接！

\textbf{证明}：设
\(\mathrm{index}(h_1) = i_1 \geq i_2 = \mathrm{index}(h_2)\)。

\begin{itemize}
\tightlist
\item
  补球面维数：\((n - i_1) - 1\)
\item
  粘接球面维数：\(i_2 - 1\)\\
\item
  维数和：\(\leq n - 2 < \dim \partial(M + h_1) = n - 1\)
\end{itemize}

由 Thom 横截性，相交维数 \(= -1\)，即不相交。\(\blacksquare\)
\end{quote}

\begin{quote}
📌 \textbf{约定}：以后总按 index 从小到大粘接。Index 相同时顺序可交换。
\end{quote}

\textbf{标准分解}：

\begin{figure}[H]

{\centering \pandocbounded{\includegraphics[keepaspectratio]{24-Smale_s_work.png}}

}

\caption{标准分解1}

\end{figure}%

\begin{figure}[H]

{\centering \pandocbounded{\includegraphics[keepaspectratio]{26-Smale_s_work.png}}

}

\caption{标准分解2}

\end{figure}%

\begin{center}\rule{0.5\linewidth}{0.5pt}\end{center}

\subsubsection{2.2.3
几何与代数消解}\label{ux51e0ux4f55ux4e0eux4ee3ux6570ux6d88ux89e3}

\begin{quote}
{[}!example{]} 🔍 例：几何消解

粘接相邻维数的 handle 时，可能出现消解：

\begin{figure}[H]

{\centering \pandocbounded{\includegraphics[keepaspectratio]{15-Smale_s_work.png}}

}

\caption{几何消解1}

\end{figure}%

在 \(M\) 上粘 0-handle \(h_1\) 和 1-handle \(h_2\)，结果收缩到 \(M\)。

🔑 \textbf{关键}：\(h_1\) 的补球面与 \(h_2\)
的粘接球面恰好\textbf{相交于 1 点}！

\begin{figure}[H]

{\centering \pandocbounded{\includegraphics[keepaspectratio]{16-Smale_s_work.png}}

}

\caption{几何消解2}

\end{figure}%
\end{quote}

\begin{quote}
{[}!note{]} 📜 Lemma 3：几何消解

对于 \(M + h_1 + h_2\)，若 \(\mathscr{S}\) 与 \(\mathcal{S}\)
\textbf{恰好横截相交于一点}，则 \[\boxed{M + h_1 + h_2 \cong M}\]

⚠️ 由 Thom 定理，仅当 \(\mathrm{index}(h_2) = \mathrm{index}(h_1) + 1\)
时发生。

\begin{center}\rule{0.5\linewidth}{0.5pt}\end{center}

\textbf{证明}：将 \(h_1, h_2\) 分割成三部分 \(A, B, C\)（都是
\(D^3\)）：

\begin{figure}[H]

{\centering \pandocbounded{\includegraphics[keepaspectratio]{18-Smale_s_work.png}}

}

\caption{分割1}

\end{figure}%

\begin{figure}[H]

{\centering \pandocbounded{\includegraphics[keepaspectratio]{19-Smale_s_work.png}}

}

\caption{分割2}

\end{figure}%

粘接 = \(((M \cup A) \cup B) \cup C\)。关注交集：

\begin{figure}[H]

{\centering \pandocbounded{\includegraphics[keepaspectratio]{21-Smale_s_work.png}}

}

\caption{交集}

\end{figure}%

沿 \(D^2\) 粘 \(D^3\) 不改变拓扑（立方体放桌面）：

\begin{figure}[H]

{\centering \pandocbounded{\includegraphics[keepaspectratio]{22-Smale_s_work.png}}

}

\caption{拓扑不变}

\end{figure}%

故 \(M + h_1 + h_2 \cong M\)。\(\blacksquare\)
\end{quote}

\begin{quote}
{[}!note{]} 📜 Lemma 4：代数消解 (Whitney Trick)

考虑\textbf{代数相交数}
\(\displaystyle\sigma(\mathcal{S}, \mathscr{S}) := \sum_{x \in \mathscr{S} \cap \mathcal{S}} \sigma(x)\)

\begin{center}\rule{0.5\linewidth}{0.5pt}\end{center}

\textbf{条件}： - \(N = \partial(M + h_1)\) 单连通 -
\(\dim M = n \geq 6\) - \(\mathrm{index}(h_1) = i\) 满足
\(2 \leq i \leq n - 3\)

\textbf{结论}：若代数相交数 \(= \pm 1\)，可同痕使几何相交数
\(= 1\)，从而 \(M + h_1 + h_2 \cong M\)。

\begin{center}\rule{0.5\linewidth}{0.5pt}\end{center}

\textbf{证明}：只需消去相邻的、符号相反的两交点 \(x, y\)：

\begin{figure}[H]

{\centering \pandocbounded{\includegraphics[keepaspectratio]{23-Smale_s_work.png}}

}

\caption{Whitney trick}

\end{figure}%

\begin{enumerate}
\def\labelenumi{\arabic{enumi}.}
\tightlist
\item
  在 \(\mathscr{S}\) 和 \(\mathcal{S}\) 上分别用 arc 连接 \(x, y\)（arc
  在 \(x, y\) 外不相交）→ 得 loop
\item
  由 \(N\) 单连通 → loop bounds \(D^2\)
\item
  由 Whitney 嵌入定理 + \(n \geq 6\) → \(D^2\) 光滑嵌入

  \begin{itemize}
  \tightlist
  \item
    \emph{嵌入说明}：\(\dim N \geq 5\)，两 disk 维数和
    \(= 4 < 5\)，故不自交
  \end{itemize}
\item
  沿 \(D^2\) 将 \(\mathscr{S}\) ``划出'' → 交点 \(-2\)，代数相交数不变
\item
  从最内层开始逐步操作 → 最终几何相交数 \(= 1\)
\item
  用 Lemma 3 完成。\(\blacksquare\)
\end{enumerate}
\end{quote}

\begin{center}\rule{0.5\linewidth}{0.5pt}\end{center}

\subsubsection{2.2.4 Handle Addition}\label{handle-addition}

\emph{限制到只有 \(i\)-handle 和 \((i+1)\)-handle 的情形。}

对于 \(M + h_1^i + h_2^i + h_3^{i+1}\)，若 \(h_1, h_2\) 满足消解条件，得
\(M + h_3\)。但 \textbf{\(h_3\) 的粘接球面如何变化}？

\begin{quote}
{[}!info{]} 📐 定义：Incidence Matrix \(D^{i+1}\)

设有 \(m_i\) 个 \(i\)-handle（补球面
\(\mathcal{S}_1, \ldots, \mathcal{S}_{m_i}\)），\(m_{i+1}\) 个
\((i+1)\)-handle（粘接球面
\(\mathscr{S}_1, \ldots, \mathscr{S}_{m_{i+1}}\)）。

\[\boxed{D^{i+1}_{j,k} := \sigma(\mathcal{S}_j, \mathscr{S}_k)}\]

得 \(m_i \times m_{i+1}\) 矩阵，蕴含 handle 分解的\textbf{全部信息}。

💡 \textbf{与同调的联系}：handle 分解 → cell 分解，微分
\(C_{i+1} \to C_i\) 恰为 \(D^{i+1}\)。
\end{quote}

\begin{quote}
{[}!note{]} 📜 Lemma 5：Handle Addition

可改变 \((i+1)\)-handle 的 attaching sphere 而\textbf{保持结果不变}：

\[\mathscr{S}_1 \leadsto \mathscr{S}_1' = \mathscr{S}_1 \mathbin{\#} \mathscr{S}_2 \implies M +_{\mathscr{S}_1} h_1 + \cdots \cong M +_{\mathscr{S}_1'} h_1 + \cdots\]

\begin{figure}[H]

{\centering \pandocbounded{\includegraphics[keepaspectratio]{30-Smale_s_work.png}}

}

\caption{Handle addition}

\end{figure}%

🔑 \textbf{对关联矩阵}：第三类基本列变换（加法/减法均可）！

\begin{center}\rule{0.5\linewidth}{0.5pt}\end{center}

\textbf{条件}：\(2 \leq i \leq n - 3\)

\textbf{证明}：

\begin{figure}[H]

{\centering \pandocbounded{\includegraphics[keepaspectratio]{27-Smale_s_work.png}}

}

\caption{证明1}

\end{figure}%

\begin{figure}[H]

{\centering \pandocbounded{\includegraphics[keepaspectratio]{28-Smale_s_work.png}}

}

\caption{证明2}

\end{figure}%

\begin{center}\rule{0.5\linewidth}{0.5pt}\end{center}

\emph{Index 限制的说明}：找 arc \(\alpha\) 需避开粘接球面和补球面。

\begin{longtable}[]{@{}ccc@{}}
\toprule\noalign{}
对象 & 维数 & 不相交条件 \\
\midrule\noalign{}
\endhead
\bottomrule\noalign{}
\endlastfoot
粘接球面 & \(i - 1\) & \(n - 1 > 1 + (i-1) \Rightarrow i \leq n-2\) \\
补球面 & \(n - i\) & \(n - 1 > (n-i) + 1 \Rightarrow i \geq 3\) \\
\end{longtable}

综合：\(3 \leq i \leq n - 2\)，即
\(2 \leq i - 1 \leq n - 3\)。\(\blacksquare\)
\end{quote}

\begin{center}\rule{0.5\linewidth}{0.5pt}\end{center}

\subsubsection{2.2.5 h-Cobordism
定理的证明}\label{h-cobordism-ux5b9aux7406ux7684ux8bc1ux660e}

\begin{quote}
{[}!note{]} 📜 Prop 2.42

\begin{figure}[H]

{\centering \pandocbounded{\includegraphics[keepaspectratio]{33-Smale_s_work.png}}

}

\caption{Prop 2.42}

\end{figure}%

💡 \textbf{动机}：Lemma 4, 5 要求 \(2 \leq i \leq n - 3\)，故 handle
index 落在 \([2, n-2]\)。需先找到 \textbf{regular form}（排除区间外的
handle）。
\end{quote}

\begin{quote}
{[}!tip{]} ✏️ h-Cobordism 定理的证明

设 \(W^n\) 是 h-cobordism，取 regular
分解：\(W = \partial^- W \times I + h_1 + \cdots + h_m\)。

\begin{center}\rule{0.5\linewidth}{0.5pt}\end{center}

\textbf{🎯 核心命题}：若 \(m > 0\)，存在 handle 更少的 regular 分解。

\begin{quote}
若成立，取 \(m\) 最小者 → \(m = 0\) → \(W = \partial^- W \times I\)。✓
\end{quote}

\begin{center}\rule{0.5\linewidth}{0.5pt}\end{center}

\textbf{命题证明}：

\textbf{Step 1️⃣} 设最小 index 为 \(i - 1\)，考虑所有 \((i-1)\)-handle 和
\(i\)-handle。

\textbf{Step 2️⃣} 考虑 \(H_{i-1}(W, \partial^- W)\)： - 由同伦等价 →
\(= 0\) - 由胞腔同调 \(H_{i-1} = C_{i-1}/\mathrm{Im}(D^i)\) → \(D^i\)
\textbf{满射}（满秩）

\textbf{Step 3️⃣} 对 \(D^i\) 做初等变换：

\[D^i \xrightarrow{\text{行列变换}} \text{上三角} \xrightarrow{\text{满射} \Rightarrow \text{对角}=\pm1} \text{对角矩阵}\]

\begin{figure}[H]

{\centering \pandocbounded{\includegraphics[keepaspectratio]{32-Smale_s_work.png}}

}

\caption{矩阵变换1}

\end{figure}%

\begin{figure}[H]

{\centering \pandocbounded{\includegraphics[keepaspectratio]{31-Smale_s_work.png}}

}

\caption{矩阵变换2}

\end{figure}%

\textbf{Step 4️⃣} 在对角矩阵上用 Lemma 3, 4 → handle
数减少。\(\blacksquare\)

\begin{center}\rule{0.5\linewidth}{0.5pt}\end{center}

⚠️ \textbf{补充命题}： - \textbf{命题 1} (Prop 2.42)：regular 分解存在 -
\textbf{命题 2}：\(\pi_1(N_{i-1}) = 1\)
\end{quote}

\begin{center}\rule{0.5\linewidth}{0.5pt}\end{center}

\begin{quote}
{[}!note{]} 📜 Lemma 2.45

\begin{figure}[H]

{\centering \pandocbounded{\includegraphics[keepaspectratio]{34-Smale_s_work.png}}

}

\caption{Lemma 2.45}

\end{figure}%

取 \(X = A - \{0\} \simeq B\)，\(Y = \mathrm{int}(D^s)\)，则
\(X \cap Y = S^{s-1}\)，\(X \cup Y = A\)。

对 \((A, A-\{0\}, \mathrm{int}(D^s))\) 用 Van Kampen 即证。
\end{quote}

\begin{quote}
{[}!tip{]} ✏️ 命题 2 的证明：\(\pi_1(N_{i-1}) = 1\)

\emph{策略}：分别将 \(N_{i-1}\) 和 \(W\) 与 \(W_{i-1}\) 联系。

\begin{center}\rule{0.5\linewidth}{0.5pt}\end{center}

\textbf{(I) \(W\) 与 \(W_{i-1}\)}

\(W = W_{i-1}\) + (index \(\geq i \geq 3\) 的 handle) →
\(\pi_1(W) \cong \pi_1(W_{i-1})\)

\begin{center}\rule{0.5\linewidth}{0.5pt}\end{center}

\textbf{(II) \(N_{i-1}\) 与 \(W_{i-1}\)}

取对偶：\(W_{i-1} = N_{i-1} \times I\) + (index \(= n - (i-1) \geq 3\)
的 handle)

→ \(\pi_1(N_{i-1}) \cong \pi_1(W_{i-1})\)

\begin{center}\rule{0.5\linewidth}{0.5pt}\end{center}

\textbf{(III)} 由 \(\pi_1(W) = 1\) →
\(\pi_1(N_{i-1}) = 1\)。\(\blacksquare\)
\end{quote}

\begin{center}\rule{0.5\linewidth}{0.5pt}\end{center}

\begin{quote}
{[}!tip{]} ✏️ 命题 1 的证明：Regular Handle 分解总存在

\textbf{策略}：证 0, 1-handle 可消解 → 对偶 → \((n-1), n\)-handle
也可消解。

\begin{center}\rule{0.5\linewidth}{0.5pt}\end{center}

\textbf{🔄 对偶论证框架}

\begin{longtable}[]{@{}clc@{}}
\toprule\noalign{}
步骤 & 操作 & 结果 \\
\midrule\noalign{}
\endhead
\bottomrule\noalign{}
\endlastfoot
1 & 消解 0, 1-handle & \(t_1\) \\
2 & 对偶 & \(t_2\) \\
3 & 消解 0, 1-handle（= \(t_1\) 的 \(n{-}1, n\)） & \(t_3\) \\
4 & 对偶 & \(t_4\) ✓ \\
\end{longtable}

\begin{center}\rule{0.5\linewidth}{0.5pt}\end{center}

\textbf{🔹 消解 0-Handle}

由 \(H_0(W, \partial^- W) = 0\)（空间连通），每个 0-handle
有\textbf{唯一} 1-handle 连接主体。

→ \(D^1\) 天然是\textbf{单位阵} → 直接消解所有 0-handle。

\begin{center}\rule{0.5\linewidth}{0.5pt}\end{center}

\textbf{🔹 消解 1-Handle (Handle Trading)}

🔑 \textbf{核心}：用 \textbf{3-handle 替换 1-handle}（总数不变）

\textbf{准备}： - 0-handle 已消解 → 1-handle 直接粘到
\(\partial^- W \times I\)，形成 loop - 在 \(W_2\) 上考虑（与 \(W\) 差
index \(\geq 3\) 的 handle，基本群同） - 由命题 2
策略：\(\partial^+ W_2\) 单连通

\begin{center}\rule{0.5\linewidth}{0.5pt}\end{center}

\textbf{构造}：

\textbf{Step 1️⃣} 任取 1-handle \(h_1\)，其 loop 在 \(\partial^+ W_2\)
上。由 \(\pi_1 = 1\)，bounds disk \(D^2\)（Thom + dim arg →
embedded）。记 \(\partial D^2 = \mathscr{S}_0\)。

\textbf{Step 2️⃣} 由 Lemma 3，引入互相消解的 2-handle \(h_2\) + 3-handle
\(h_3\)，使 \(W_2 = W_2 + h_2 + h_3\)。要求 \(h_2\) 避开 \(D^2\) 和其他
2-handle。记粘接球面 \(\mathscr{S}_2\)。

\textbf{Step 3️⃣} 取
\(\mathscr{S}_2' := \mathscr{S}_2 \mathbin{\#} \mathscr{S}_0\)，对应
\(h_2'\)。

\begin{figure}[H]

{\centering \pandocbounded{\includegraphics[keepaspectratio]{37-Smale_s_work.png}}

}

\caption{Handle trading}

\end{figure}%

\begin{center}\rule{0.5\linewidth}{0.5pt}\end{center}

\textbf{🔑 关键观察}： - \(\mathscr{S}_0\) bounds disk →
\(\mathscr{S}_2' \sim \mathscr{S}_2\) - \(\mathscr{S}_2'\) 与 \(h_1\)
的补球面\textbf{几何相交数 = 1}

\begin{center}\rule{0.5\linewidth}{0.5pt}\end{center}

\textbf{消解计算}：

\[\begin{aligned}
W_2 &= W_2 + h_2 + h_3 \\[4pt]
    &= W_2 + h_2' + h_3 \\[4pt]
    &= \partial^- W \times I + \cdots + \underbrace{(h_1 + h_2')}_{\text{消解}} + \cdots + h_3 \\[4pt]
    &= \partial^- W \times I + \cdots + h_3
\end{aligned}\]

\begin{center}\rule{0.5\linewidth}{0.5pt}\end{center}

\[\boxed{\text{1-handle } h_1 \leadsto \text{3-handle } h_3} \quad \blacksquare\]
\end{quote}




\end{document}
